\section{Limitations and Future Scope}

The platform has several documented limitations. The Riskometer is present for approximately 70 percent of merged fund rows; the remaining rows, dominated by inactive legacy schemes without details disclosures, are excluded by the engine rather than imputed. NAV precision on certain liquid and overnight funds can inflate volatility estimates due to rounding of rupee-denominated NAV steps; this artifact is documented rather than smoothed. Portfolio holdings are available for only 24 funds, so concentration and allocation metrics serve as enrichment context and are never scored. The repository holds no verified user-fund interaction dataset, so the recommendation methodology remains deterministic and rule-based rather than collaborative-filtering or machine-learning driven. The recommendation engine depends on the availability and quality of historical NAV data, and funds with insufficient history are excluded or receive neutral scores rather than estimates.

The assessment and Digital Twin engines do not model investment returns, interest, inflation, taxes, salary growth, future income changes, borrowing or asset sales. Goal projections assume constant monthly contributions and no returns. The one-time unexpected expense is treated as a single shock to accessible savings rather than a recurring or amortized expense. Heuristic thresholds, such as the three-month emergency coverage minimum and the 40 percent EMI burden limit, are educational rules of thumb rather than regulatory standards or calibrated risk measures.

Future work includes extending Riskometer coverage as SEBI disclosure requirements evolve, improving NAV precision handling for low- denomination funds, expanding portfolio holdings coverage as more AMCs disclose, collecting user-fund interaction data to evaluate content-based learning extensions against a majority-class baseline with documented leakage controls, and applying the Supabase migrations in a hosted environment for end-to-end authentication and persistence verification. Longer-term validation would require independent expert review of recommendation appropriateness, user comprehension studies with Indian beginners across languages and regions, and longitudinal or randomized studies of decision support and outcomes, none of which have been performed.